\subsection{技术指标的真实角色}

\begin{frame}
    \frametitle{技术指标的问题}

    \alert{不能将技术指标当做买卖信号} \\
    \hspace{2em} 指标是滞后的：MACD 金叉时，价格已经涨了一段 \\
    \hspace{2em} 指标是简化的：一个数字无法概括复杂的市场状态 \\
    \hspace{2em} 指标是被利用的：当所有人都在看同一个信号，信号就失效了

    ~\par
    指标不是用来直接交易的，而是用来描述市场状态的
\end{frame}

\begin{frame}
    \frametitle{指标的本质=数学变换}
    
    \alert{数学变换}: 指标不会创造新信息，只是换一种方式呈现已有信息

    \alert{滞后性}: 任何基于历史数据计算的指标，都必然是滞后的

    \alert{信息压缩与损失}: 将多个交易日的数据，压缩为一个数字，损失大量信息

    \begin{table}[htbp]
        \centering
        % \caption{常见技术指标的构造方式}
        \begin{tabular}{lll}
            \toprule
            原始数据 & 变换方式 & 得到的指标 \\
            \midrule
            收盘价序列 & 滑动平均 & MA（移动平均线） \\
            收盘价序列 & 指数加权平均 & EMA（指数移动平均） \\
            涨跌幅序列 & 相对强度计算 & RSI \\
            价格序列 & 均值 $\pm$ 标准差 & 布林带 \\
            最高价、最低价、收盘价 & 多重计算 & MACD \\
            \bottomrule
        \end{tabular}
    \end{table}
\end{frame}

\begin{frame}
    \frametitle{MA / EMA: 平滑与滞后的权衡}

    \alert{移动平均线(MA)}: 过去 $N$ 天的收盘价的简单平均 \\
    \hspace{2em} $MA(5) = (P_1 + P_2 + P_3 + P_4 + P_5) / 5$ \\
    \hspace{2em} 权重分布平均，对新数据反应慢，更平滑，适合长期趋势

    ~\par
    \alert{指数移动平均线(EMA)}: 给近期价格更高的权重 \\
    \hspace{2em} $EMA = $ 今日价格 $\times \alpha$ + 昨日 $EMA \times (1 - \alpha)$, 其中 $\alpha = 2 / (N + 1)$ \\
    \hspace{2em} 权重分布近期更高，对新数据反应快，更灵敏，适合短期趋势
\end{frame}

\begin{frame}
    \frametitle{MACD: 趋势与动量}

    MACD 是最常用的趋势指标之一，由三部分组成: \\
    \hspace{2em} \alert{DIFF(快线)} $= EMA_{12} - EMA_{26}$, 短期趋势与长期趋势的差距 \\
    \hspace{2em} \alert{DEA(信号线)} $= EMA_9(DIFF)$, DIFF 的平滑版本\\
    \hspace{2em} 柱状图 = DIFF - DEA, 趋势加速/减速的程度 \\

    \begin{table}[htbp]
        \centering
        \caption{MACD 状态解读}
        \begin{tabular}{ccc}
            \toprule
            MACD 状态 & 含义 & 市场解读 \\
            \midrule
            DIFF $>$ 0, 柱状图变大 & 上涨趋势加速 & 多头强势 \\
            DIFF $>$ 0, 柱状图变小 & 上涨趋势减速 & 可能见顶 \\
            DIFF $<$ 0, 柱状图变小 & 下跌趋势减速 & 可能见底 \\
            DIFF $<$ 0, 柱状图变大 & 下跌趋势加速 & 空头强势 \\
            \bottomrule
        \end{tabular}
    \end{table}
\end{frame}

\begin{frame}
    \frametitle{背离分析}

    \alert{背离}是技术分析中最重要的概念之一: 

    \begin{table}[htbp]
        \centering
        % \caption{MACD 背离类型对照表}
        \begin{tabular}{cccc}
            \toprule
            背离类型 & 价格表现 & MACD 表现 & 含义 \\
            \midrule
            顶背离 & 创新高 & 未创新高 & 上涨动能衰竭，可能下跌 \\
            底背离 & 创新低 & 未创新低 & 下跌动能衰竭，可能上涨 \\
            \bottomrule
        \end{tabular}
    \end{table}
\end{frame}